\section{Limitations}
\label{sec:limitations}

\paragraph{Scale.} Our models have 3.2M parameters and are trained from scratch on
synthetic sequence tasks. Nothing in our argument requires scale, but nothing in
our experiments establishes that the same conditioning holds for pretrained
language models, where positional codes interact with token distributions and
with training that already spans many lengths. We regard the present study as a
controlled test of a mechanism, not as a claim about frontier models.

\paragraph{One training length.} All runs use $L_{\mathrm{tr}}{=}8$. A stronger
version of our claim would show that the coupling classification predicts the
\emph{shape} of the extrapolation curve across several training lengths, not just
the sign of a difference at one. This is the most important extension we have not
run.

\paragraph{Online sampling.} Each step draws a fresh batch rather than sampling
from a fixed finite training set. This removes memorisation as an explanation and
equalises the number of examples across conditions, but it also makes the task
easier than the fixed-train-set protocol used in some prior work, so absolute
accuracies are not directly comparable across papers.

\paragraph{Coverage of the taxonomy.} Five tasks is a small battery, and the
position-free class is represented by a single task (sort). The two coupled
tasks, while both length-dependent, differ in how the dependence enters
(complement versus rotation), so we cannot separate ``coupled'' from ``coupled in
this particular way''.

\paragraph{Confound in the learned-APE condition.} At test lengths the learned
embedding table is queried at positions never updated during training. This is
faithful to how learned absolute codes behave in deployment, but it means the
learned condition conflates ``absolute position is unhelpful'' with ``the
untrained tail of a lookup table is unhelpful''. A variant with a
length-extrapolable initialisation would separate the two; we did not run one.

\paragraph{Statistical power.} Two seeds per cell leave only 25 residual degrees
of freedom for the interaction test, and treating seeds as replicates is an
approximation. The interaction we report is far from marginal
($F(16,25)=32.3$), but the cell-level means it is computed from are noisy enough
that we would not defend a claim about any single pairwise difference. A
five-seed replication is the obvious next step and we have not run it.

\paragraph{Behavioural evidence only.} We report accuracy, not mechanism. Claims
about \emph{why} a coupled task is harder---whether models recover length from
attention statistics, whether they implement a position-arithmetic circuit---would
require probing or circuit analysis that we have not performed.
